%%%%%%%%%%%%%%%%%%%%%%%%%%%%%%%%%%%%%%%%%%%%%%%%%%%%%%%%%%%%%%%%%%%%
%% I, the copyright holder of this work, release this work into the
%% public domain. This applies worldwide. In some countries this may
%% not be legally possible; if so: I grant anyone the right to use
%% this work for any purpose, without any conditions, unless such
%% conditions are required by law.
%%%%%%%%%%%%%%%%%%%%%%%%%%%%%%%%%%%%%%%%%%%%%%%%%%%%%%%%%%%%%%%%%%%%

\documentclass{beamer}
\usetheme[faculty=ped,basePath=../fibeamer,logo=../fibeamer/logo/mu/Logo-UC-3.png]{fibeamer}
\graphicspath{{../resources/}}

\usepackage{algorithm}
\usepackage{algpseudocode}

\usetheme[faculty=ped]{fibeamer}
\usepackage[utf8]{inputenc}
\usepackage[
  main=english, %% By using `czech` or `slovak` as the main locale
                %% instead of `english`, you can typeset the
                %% presentation in either Czech or Slovak,
                %% respectively.
  czech, slovak %% The additional keys allow foreign texts to be
]{babel}        %% typeset as follows:
%%
%%   \begin{otherlanguage}{czech}   ... \end{otherlanguage}
%%   \begin{otherlanguage}{slovak}  ... \end{otherlanguage}
%%
%% These macros specify information about the presentation
\title{Control de flujo} %% that will be typeset on the
\subtitle{Semestre II-2026} %% title page.
\author{Andreina Cota Vidaurre}
%% These additional packages are used within the document:
\usepackage{ragged2e}  % `\justifying` text
\usepackage{booktabs}  % Tables
\usepackage{tabularx}
\usepackage{tikz}      % Diagrams
\usetikzlibrary{calc, shapes, backgrounds}
\usetikzlibrary{positioning, arrows.meta, shadows.blur, shapes.geometric}
\usepackage{amsmath, amssymb}
\usepackage{url}       % `\url`s
\usepackage{listings}  % Code listings
\usepackage{graphicx}
\usepackage{amsmath}
\usepackage[T1]{fontenc}
\usepackage{xcolor}
\usepackage{minted}
% \usepackage{adjustbox}


% Definicion de pseudocodigo en espaniol
\lstdefinelanguage{PseudocodigoES}{
    morekeywords={INICIO,FIN,PARA,HACER,SI,ENTONCES,FIN_SI,FIN_PARA,IMPRIMIR},
    sensitive=true,
    morecomment=[l]{//},
    morestring=[b]"
}

% Definicion de pseudocodigo en ingles
\lstdefinelanguage{PseudocodeEN}{
    morekeywords={BEGIN,END,FOR,DO,IF,THEN,END_IF,END_FOR,PRINT},
    sensitive=true,
    morecomment=[l]{//},
    morestring=[b]"
}

\lstset{
    basicstyle=\ttfamily\small,
    keywordstyle=\bfseries,
    columns=fullflexible,
    keepspaces=true,
    showstringspaces=false,
    frame=single,
    xleftmargin=0em,
    framexleftmargin=0em,
    framesep=2pt,
    gobble=4,
    resetmargins=true,
    aboveskip=0pt,
    belowskip=0pt
}

\lstdefinestyle{pythonstyle}{
    language=Python,
    basicstyle=\ttfamily\small,
    keywordstyle=\bfseries,
    commentstyle=\itshape,
    stringstyle=\ttfamily,
    showstringspaces=false,
    columns=fullflexible,
    keepspaces=true,
    frame=single,
    breaklines=true,
    xleftmargin=0em,
    framexleftmargin=0em,
    framesep=2pt,
    aboveskip=0pt,
    belowskip=0pt,
    gobble=4,
    resetmargins=true,
    emph={print,input,int,float,str,bool,len,range},
    emphstyle=\bfseries
}

\lstset{style=pythonstyle}

% Estilos del diagrama
\tikzstyle{iniciofin} = [
    ellipse,
    minimum width=2.5cm,
    minimum height=1cm,
    text centered,
    draw=black,
    fill=blue!25
]

\tikzstyle{proceso} = [
    rectangle,
    rounded corners,
    minimum width=3cm,
    minimum height=1cm,
    text centered,
    draw=black,
    fill=cyan!20
]

\tikzstyle{decision} = [
    diamond,
    aspect=2,
    text centered,
    draw=black,
    fill=yellow!35
]

\tikzstyle{flecha} = [
    thick,
    ->,
    >=Stealth
]

% \definecolor{green_python}{HTML}{52A736}

% \newcolumntype{Y}{>{\raggedright\arraybackslash}X}

\frenchspacing
\begin{document}
  \shorthandoff{-}
  \frame[c]{\maketitle}

% \begin{darkframes}

    \begin{frame}[fragile,t]{Contenido}
        \begin{enumerate}
            \item Control de flujo
            \item Estructura \mintinline{python}{if}
            \item Estructura \mintinline{python}{if - else}
            \item Estructura \mintinline{python}{if - elif - else}
            \item Estructura \mintinline{python}{if - elif}
        \end{enumerate}
    \end{frame}

    \begin{frame}{Control de flujo}
        \begin{itemize}
            \item Es el orden en el que se ejecutan las instrucciones de un programa
            \item Por defecto, el código se ejecuta secuencialmente (línea por línea, de arriba hacia abajo)
            \item Permite que, según el resultado, el programa continue por un camino u otro
            \item El programa no siempre ejecuta lo mismo
            \item Puede cambiar su camino según la condición
        \end{itemize}
    \end{frame}

    \begin{frame}{Tipos de estructuras de control}
        \begin{itemize}
            \item Secuenciales: instrucciones que se ejecutan una después de la otra
            \item Condicionales (selección): permiten tomar decisiones (\mintinline{python}{if}, \mintinline{python}{else}, \mintinline{python}{elif})
            \item Iterativas (repetición): permiten repetir bloques de código (\mintinline{python}{for}, \mintinline{python}{while})
        \end{itemize}
    \end{frame}

    \begin{frame}{Condición}
        \begin{itemize}
            \item Una condición permite al programa \textbf{tomar decisiones}
            \item Una condición evalua una expresión como \mintinline{python}{True o False}
            \item Para formar condiciones se usan comparadores como: \mintinline{python}{>, <, ==, !=}
            \item Las condiciones suelen trabajar con variables, números o textos
            \item El resultado de la condición define qué instrucciones se ejecutan
        \end{itemize}
    \end{frame}

    \begin{frame}[fragile,t]{Condición}
        \begin{minted}[xleftmargin=-2cm, frame=leftline, fontsize=\small]{python}
            edad >= 18                  # mayor de edad o no
            nota >= 4                   # aprobado o reprobado
            temperatura >= 37.5         # tiene fiebre o no
            password == "123456"        # contraseña correcta o no
            aprobado != True            # no aprobado o no
        \end{minted}
    \end{frame}

    \begin{frame}[fragile,t]{\mintinline{python}{if}}
         \begin{itemize}
            \item Se usa para tomar una decisión en el programa
            \item Evalua una condición
            \item Si la condición es \mintinline{python}{True}, ejecuta un bloque de código
            \item Si la condición es \mintinline{python}{False}, no ejecuta ese bloque
            \item En Python, se escribe con \mintinline{python}{:} e indentación
            \item En Python, el bloque se reconoce por la indentación
         \end{itemize}
    \end{frame}

    \begin{frame}{Diagrama de flujo: \mintinline{python}{if}}
        \centering
        \begin{tikzpicture}[node distance=0.8cm, 
            every node/.style={scale=0.75}, 
            transform shape]
        
            % Nodos
            \node (inicio) [iniciofin] {Inicio};
            \node (leer) [proceso, below=of inicio] {Leer combustible};
            
            \node (condicion) [decision, below=0.7cm of leer] 
                {combustible >= 50?};
            
            \node (salida) [proceso, below=0.7cm of condicion] 
                {Autorizar salida};
            
            \node (fin) [iniciofin, below=0.7cm of salida] {Fin};
        
            % Flechas
            \draw [flecha] (inicio) -- (leer);
            \draw [flecha] (leer) -- (condicion);
        
            % Si
            \draw [flecha] (condicion) -- 
                node[anchor=east] {Si} (salida);
        
            \draw [flecha] (salida) -- (fin);
        
            % No: omite el bloque del if
            \draw [flecha] (condicion.east) -- ++(2,0)
                node[midway, above] {No}
                |- (fin.east);
        
        \end{tikzpicture}
    \end{frame}


    \begin{frame}[fragile,t]{\mintinline{python}{if} - Ejemplos}
        1. Comprobar si un estudiante aprobó la prueba de programación
        \begin{minted}[xleftmargin=-2cm, frame=leftline, fontsize=\small]{python}
            puntaje_interrogacion = 5.4
            if(puntaje_interrogacion >= 4):
                print("El estudiante aprobo la prueba")
        \end{minted}
        2. Se necesita comprobar si el voltaje medido en un circuito es igual al voltaje nominal 
        esperado (220V). Si coincide, se confirma que el circuito funciona correctamente.
        \begin{minted}[xleftmargin=-2cm, frame=leftline, fontsize=\small]{python}
            voltaje_medido = 225
            if voltaje_medido == 220:
                print("El circuito funciona correctamente")
        \end{minted}
    \end{frame}

    \begin{frame}[fragile,t]{\mintinline{python}{if} - Ejemplos}
        3. En un sistema de seguridad, hay una variable que indica si la puerta está 
        abierta (\mintinline{python}{True}) o cerrada (\mintinline{python}{False}). 
        Si la puerta está abierta, se debe mostrar una advertencia.
        \begin{minted}[xleftmargin=-2cm, frame=leftline, fontsize=\small]{python}
            puerta_abierta = True
            if puerta_abierta:
                print("Advertencia: La puerta está abierta")
        \end{minted}
        4. Un sistema recibe el estado de una máquina como texto. Si el estado es "detenida", 
        se debe mostrar un mensaje indicando que la máquina no está operando.
        \begin{minted}[xleftmargin=-2cm, frame=leftline, fontsize=\small]{python}
            estado_maquina = "detenida"
            if estado_maquina == "detenida":
                print("La máquina no está operando")
        \end{minted}
    \end{frame}

    % \begin{frame}{\mintinline{python}{if} - Actividad}
    %     Reconocer variables y construir condiciones
    %     \begin{itemize}
    %         \item Un soldado aprobó la práctica de tiro si su puntaje es mayor o igual a 70
    %         \item Un militar llegó puntual a la formación si su hora de llegada es menor o igual a 6
    %         \item Un recluta pasa a la siguiente fase de entrenamiento si su nota final es mayor o igual a 80
    %         \item Se debe mostrar un mensaje si la cantidad de munición disponible es diferente de 0
    %     \end{itemize}
    % \end{frame}

    \begin{frame}[fragile,t]{\mintinline{python}{if - else}}
        \begin{itemize}
            \item \mintinline{python}{if - else} permite al programa tomar una decisión entre dos caminos
            \item Primero se evalua una condición
            \item Si la condición es \mintinline{python}{True}, se ejecuta el bloque de \mintinline{python}{if}
            \item Si la condición es \mintinline{python}{False}, se ejecuta el bloque de \mintinline{python}{else}
            \item En Python, ambos bloques se definen con \mintinline{python}{:} e indentación
        \end{itemize}
    \end{frame}

    \begin{frame}{Diagrama de flujo: \mintinline{python}{if - else}}
        \centering
        \begin{tikzpicture}[node distance=1.4cm, 
        every node/.style={scale=0.75}, 
        transform shape]
        
            % Nodos
            \node (inicio) [iniciofin] {Inicio};
            \node (leer) [proceso, below of=inicio] {Leer edad};
            \node (condicion) [decision, below of=leer, yshift=-0.3cm] {edad >= 18?};
            \node (mayor) [proceso, below left=1.4cm and 2cm of condicion] {Mostrar: Mayor de edad};
            \node (menor) [proceso, below right=1.4cm and 2cm of condicion] {Mostrar: Menor de edad};
            \node (fin) [iniciofin, below=2cm of condicion] {Fin};
        
            % Flechas
            \draw [flecha] (inicio) -- (leer);
            \draw [flecha] (leer) -- (condicion);
            \draw [flecha] (condicion) -- node[anchor=east] {Si} (mayor);
            \draw [flecha] (condicion) -- node[anchor=west] {No} (menor);
            \draw [flecha] (mayor) |- (fin);
            \draw [flecha] (menor) |- (fin);
        
        \end{tikzpicture}
    \end{frame}

    \begin{frame}[fragile,t]{\mintinline{python}{if} - Ejemplos}
        1. Comprobar si un estudiante aprobó la prueba de programación. 
        \textbf{Si no aprobó, se debe mostrar un mensaje de reprobación.}
        
        \begin{minted}[xleftmargin=-2cm, frame=leftline, fontsize=\small]{python}
            puntaje_interrogacion = 5.4
            if(puntaje_interrogacion >= 4):
                print("El estudiante aprobo la prueba")
            else:
                print("El estudiante no aprobo la prueba")
        \end{minted}
    \end{frame}

    \begin{frame}[fragile,t]{\mintinline{python}{if - else} - Ejemplos}
        2. Se necesita comprobar si el voltaje medido en un circuito es igual al voltaje nominal 
        esperado (220V). Si coincide, se confirma que el circuito funciona correctamente. \textbf{Si no, 
        se debe mostrar un mensaje de error.}
        \begin{minted}[xleftmargin=-2cm, frame=leftline, fontsize=\small]{python}
            voltaje_medido = 225
            if voltaje_medido == 220:
                print("El circuito funciona correctamente")
            else:
                print("El circuito no funciona correctamente")
        \end{minted}
    \end{frame}

    \begin{frame}[fragile,t]{\mintinline{python}{if - else} - Ejemplos}
        3. En un sistema de seguridad, hay una variable que indica si la puerta está 
        abierta (\mintinline{python}{True}) o cerrada (\mintinline{python}{False}). 
        Si la puerta está abierta, se debe mostrar una advertencia. \textbf{Si está cerrada, se debe 
        mostrar un mensaje de confirmación.}
        \begin{minted}[xleftmargin=-2cm, frame=leftline, fontsize=\small]{python}
            puerta_abierta = True
            if puerta_abierta:
                print("Advertencia: La puerta está abierta")
            else:
                print("La puerta está cerrada")
        \end{minted}
    \end{frame}

    \begin{frame}[fragile,t]{\mintinline{python}{if - else} - Ejemplos}
        4. Un sistema recibe el estado de una máquina como texto. Si el estado es "detenida", 
        se debe mostrar un mensaje indicando que la máquina no está operando. \textbf{Si está operando, 
        se debe mostrar un mensaje de confirmación.}
        \begin{minted}[xleftmargin=-2cm, frame=leftline, fontsize=\small]{python}
            estado_maquina = "detenida"
            if estado_maquina == "detenida":
                print("La máquina no está operando")
            else:
                print("La máquina está operando")
        \end{minted}
    \end{frame}

    % \begin{frame}[fragile,t]{\mintinline{python}{if - else} - Ejemplos}
    %     1. Ingreso al campo de instrucción
    %     \begin{minted}[xleftmargin=-2cm, frame=leftline, fontsize=\small]{python}
    %         edad = 17
    %         if edad >= 18:
    %             print("El cadete puede ingresar al campo ...")
    %         else:
    %             print("El cadete no puede ingresar al campo ...")
    %     \end{minted}
    %     2. Aprobación de prueba física
    %     \begin{minted}[xleftmargin=-2cm, frame=leftline, fontsize=\small]{python}
    %         puntaje_fisico = 58
    %         if puntaje_fisico >= 60:
    %             print("El soldado aprobo la prueba fisica")
    %         else:
    %             print("El soldado no aprobo la prueba fisica")
    %     \end{minted}
    % \end{frame}

    % \begin{frame}[fragile,t]{\mintinline{python}{if - else} - Actividad}
    %     Reconoce las variables y construye las condiciones
    %     \begin{itemize}
    %         \item Un militar llego puntual si su hora de llegada es menor o igual a 6. Si llega después, se considera tarde
    %         \item Si la cantidad de munición es mayor que 0, la unidad puede iniciar el entrenamiento. Si no hay munición, el entrenamiento no puede comenzar
    %         \item Si un recluta asistió a las 5 sesiones programadas, puede rendir la evaluación final. Si no asistió a todas, no puede rendirla
    %         \item Si la distancia entre dos unidades es mayor o igual a 100 metros, se considera segura. Si es menor, se debe emitir una advertencia.
    %     \end{itemize}
    % \end{frame}

    \begin{frame}[fragile,t]{\mintinline{python}{if - elif - else}}
        \begin{itemize}
            \item Permite tomar decisiones entre varias opciones
            \item Las condiciones se evaluan una por una y en orden
            \item \mintinline{python}{if} inicia la primera condición
            \item \mintinline{python}{elif} agrega condiciones adicionales
            \item \mintinline{python}{else} representa el caso final cuando ninguna se cumple
        \end{itemize}
        El orden importa. \textbf{Cuando una condición se cumple, las demas ya no se revisan.}
    \end{frame}

    \begin{frame}{Diagrama de flujo: \mintinline{python}{if - elif - else}}
        \centering
        \begin{tikzpicture}[node distance=0.5cm, 
            every node/.style={scale=0.72}, 
            transform shape]
        
            % Nodos
            \node (inicio) [iniciofin] {Inicio};
            \node (leer) [proceso, below=of inicio] {Leer nivel de alerta};
        
            % Primera condicion
            \node (condicion1) [decision, below=0.5cm of leer] 
                {nivel >= 3?};
        
            % Segunda condicion
            \node (condicion2) [decision, below=0.5cm of condicion1] 
                {nivel == 2?};
        
            % Acciones
            \node (accion1) [proceso, right=2.5cm of condicion1] 
                {Activar protocolo};
        
            \node (accion2) [proceso, right=2.5cm of condicion2] 
                {Reforzar vigilancia};
        
            \node (accion3) [proceso, below=0.5cm of condicion2] 
                {Vigilancia normal};
        
            \node (fin) [iniciofin, below=0.5cm of accion3] {Fin};
        
            % Flechas principales
            \draw [flecha] (inicio) -- (leer);
            \draw [flecha] (leer) -- (condicion1);
        
            % Primera condicion
            \draw [flecha] (condicion1) -- 
                node[above] {Si} (accion1);
        
            \draw [flecha] (condicion1) -- 
                node[anchor=east] {No} (condicion2);
        
            % Segunda condicion
            \draw [flecha] (condicion2) -- 
                node[above] {Si} (accion2);
        
            \draw [flecha] (condicion2) -- 
                node[anchor=east] {No} (accion3);
        
            % Salidas hacia fin
            \draw [flecha] (accion1.east) -- ++(1,0) |- (fin.east);
            \draw [flecha] (accion2.east) -- ++(0.7,0) |- (fin.east);
            \draw [flecha] (accion3) -- (fin);
        
        \end{tikzpicture}
    \end{frame}

    \begin{frame}[fragile,t]{\mintinline{python}{if - elif - else} - Ejemplos}
        1. Según la temperatura del día, una persona decide qué ropa usar. 
        Si la temperatura es mayor a 28°C, usa ropa ligera. Si está entre 15°C y 28°C, 
        usa ropa normal. Si es menor a 15°C, usa abrigo.
        
        \begin{minted}[xleftmargin=-2cm, frame=leftline, fontsize=\small]{python}
            temperatura = 10

            if temperatura > 28:
                print("Usa ropa ligera")
            elif temperatura >= 15:
                print("Usa ropa normal")
            else:
                print("Usa abrigo")
        \end{minted}
    \end{frame}

    \begin{frame}[fragile,t]{\mintinline{python}{if - elif - else} - Ejemplos}
        2. Un semáforo indica qué acción debe tomar un conductor según su color. 
        Si está en rojo, debe detenerse. Si está en amarillo, debe prepararse para detenerse. 
        Si está en verde, puede avanzar.
        \begin{minted}[xleftmargin=-2cm, frame=leftline, fontsize=\small]{python}
            color = "amarillo"

            if color == "rojo":
                print("Detente")
            elif color == "amarillo":
                print("Prepárate para detenerte")
            else:
                print("Puedes avanzar")
        \end{minted}
    \end{frame}

    \begin{frame}[fragile,t]{\mintinline{python}{if - elif}}
        \begin{itemize}
            \item Permite evaluar múltiples condiciones en orden, 
            cuando hay más de dos posibles resultados
            \item Python evalúa las condiciones de arriba hacia abajo y 
            ejecuta solo el primer bloque cuya condición sea verdadera
            \item \mintinline{python}{elif} es la forma corta de "else if";
            evita tener que anidar varios \mintinline{python}{if} por separado
            \item El \mintinline{python}{else} al final es opcional, se usa para
            manejar el caso cuando ninguna de las condiciones anteriores es verdadera
            \item Se pueden usar tantos \mintinline{python}{elif} como se necesiten.
        \end{itemize}
    \end{frame}

    \begin{frame}[fragile,t]{\mintinline{python}{if - elif} - Ejemplos}
        1. Una tienda aplica descuentos según la cantidad de productos que compra un cliente. 
        Si compra más de 10 productos, obtiene 20\% de descuento. 
        Si compra entre 5 y 10 productos, obtiene 10\% de descuento.
        
        \begin{minted}[xleftmargin=-2cm, frame=leftline, fontsize=\small]{python}
            cantidad = 7

            if cantidad > 10:
                print("Descuento del 20%")
            elif cantidad >= 5:
                print("Descuento del 10%")
        \end{minted}
    \end{frame}

    \begin{frame}[fragile,t]{\mintinline{python}{if - elif} - Ejemplos}
        2. Un celular muestra un mensaje según el nivel de batería. Si la batería es menor al 15\%, 
        muestra un aviso de batería crítica. Si está entre 15\% y 30\%, 
        muestra un aviso de batería baja.
        \begin{minted}[xleftmargin=-2cm, frame=leftline, fontsize=\small]{python}
            bateria = 20

            if bateria < 15:
                print("Batería crítica")
            elif bateria <= 30:
                print("Batería baja")
        \end{minted}
    \end{frame}

    % \begin{frame}[fragile,t]{Errores comunes en condicionales}
    %     \begin{itemize}
    %         \item Confundir \mintinline{python}{=} con \mintinline{python}{==}
    %     \end{itemize}
    % \end{frame}

    % \begin{frame}[fragile,t]{Errores comunes en condicionales}
    %     Correcto:
    %     \begin{minted}[xleftmargin=-2cm, frameleftmargin=0pt]{python}
    %         if edad == 18:
    %     \end{minted}
    %     Incorrecto:
    %     \begin{minted}[xleftmargin=-2cm, frameleftmargin=0pt]{python}
    %         if edad = 18:
    %     \end{minted}
    % \end{frame}

    % \begin{frame}[fragile,t]{Errores comunes en condicionales}
    %     \begin{itemize}
    %         \item Confundir \mintinline{python}{=} con \mintinline{python}{==}
    %         \item Mala indentación
    %     \end{itemize}
    % \end{frame}

    % \begin{frame}[fragile,t]{Errores comunes en condicionales}
    %     \begin{minted}[xleftmargin=-2cm, frameleftmargin=0pt]{python}
    %         if edad >= 18:
    %         print("Mayor de edad")
    %     \end{minted}
    % \end{frame}

    \begin{frame}[fragile,t]{Errores comunes en condicionales}
        \begin{itemize}
            \item Confundir \mintinline{python}{=} con \mintinline{python}{==}
            \item Mala indentación
            \item Condiciones mal ordenadas en \mintinline{python}{elif}
            \item Escribir varios \mintinline{python}{if} en lugar de \mintinline{python}{elif}. Esto hace que Python
            evalúe todas las condiciones, incluso si ya se cumplió una anteriormente.
            \begin{minted}[xleftmargin=-2cm, frame=leftline, fontsize=\small]{python}
            if nota >= 90:
                print("A")
            if nota >= 70:
                print("B")
            \end{minted}
        \end{itemize}
    \end{frame}

    \begin{frame}[fragile,t]{Errores comunes en condicionales}
        \begin{itemize}
            \item Colocar condiciones más generales antes que las específicas.
            \begin{minted}[xleftmargin=-2cm, frame=leftline, fontsize=\small]{python}
            if nota >= 60:
                print("Aprobado")
            elif nota >= 90:
                print("Excelente")
            \end{minted}
            \item Comparar tipos de datos diferentes.
            \begin{minted}[xleftmargin=-2cm, frame=leftline, fontsize=\small]{python}
            if nota == "90":
                print("Excelente")
            \end{minted}
            \item Olvidar que \mintinline{python}{elif} y \mintinline{python}{else} 
            dependen del \mintinline{python}{if} anterior.
            \item Comparar booleanos innecesariamente: \mintinline{python}{if variable == True:}
            en lugar de \mintinline{python}{if variable:}. 
            
            No es un error, pero es una mala práctica.
        \end{itemize}
    \end{frame}

    % \begin{frame}[fragile,t][fragile,t]{Errores comunes en condicionales}
    %     \begin{minted}[xleftmargin=-2cm, frame=leftline, fontsize=\small]{python}
    %         nota = 6.7
    %         if nota >= 5:
    %             print("Aprobado")
    %         elif nota >= 6.5:
    %             print("Excelente")
    %     \end{minted}
    % \end{frame}

   
    % \begin{frame}{Actividad}
    %     \begin{enumerate}
    %         \item Pedir un número e indicar si es positivo, negativo o cero
    %         \item Pedir dos números e indicar si el primero es divisible por el segundo
    %         \item Pedir dos números y mostrar cuál es el mayor
    %         \item Dadas tres notas, mostrar la mayor y la menor
    %         \item Dadas tres calificaciones, calcular el promedio y mostrarlo redondeado a dos cifras. Luego indicar si el estudiante aprueba o no
    %         \item Dados dos valores, indicar si la diferencia entre ellos es pequeña, media o grande. La diferencia es pequeña si es menor que 5, es media si está entre 5 y 10, y es grande si es mayor que 10
    %     \end{enumerate}
    % \end{frame}

    % \begin{frame}{Actividad}
    %     \begin{enumerate}
    %         \item Pedir un número de tres cifras e indicar si es capicúa o no
    %         \item Pedir un número de tres cifras y mostrar cuál es su cifra mayor
    %         \item Pedir un número de cuatro cifras, sumar sus cifras e indicar si la suma es par o impar
    %         \item Tres escuadras tienen cuatro cantidades distintas de munición. Mostrar cuál tiene más y cuál tiene menos y si la diferencia entre ambas supera el límite de 100
    %         \item Dadas tres horas de llegada, identificar la más temprana, la más tardía y decidir si la diferencia entre ambas es aceptable. Es aceptable si la diferencia es menor a 15 minutos
    %     \end{enumerate}
    % \end{frame}

    \begin{frame}[fragile,t]{Resumen}
        \begin{itemize}
            \item El control de flujo usa el resultado de la condición para decidir qué instrucciones ejecutar
            \item Las condiciones permiten evaluar una situación
            \item Una condición siempre se evalua como \mintinline{python}{True} o \mintinline{python}{False}
            \item \mintinline{python}{if} ejecuta un bloque si la condición es verdadera
            \item \mintinline{python}{else} permite un camino alternativo
            \item \mintinline{python}{elif} permite evaluar varias posibilidades
            \item La indentación es fundamental
        \end{itemize}
    \end{frame}
    
\end{document}
